\documentclass[11pt,letterpaper]{article}
\usepackage[margin=1in]{geometry}
\usepackage[T1]{fontenc}
\usepackage{lmodern}
\usepackage{microtype}
\usepackage{amsmath,amssymb}
\usepackage{booktabs,array,enumitem,xcolor}
\usepackage{titlesec}
\usepackage[colorlinks=true,linkcolor=black,urlcolor=blue!50!black,citecolor=black]{hyperref}
\setlist{itemsep=2pt,topsep=3pt}
\setlength{\parindent}{0pt}\setlength{\parskip}{6pt}
\titleformat{\section}{\normalfont\large\bfseries}{\thesection}{0.6em}{}
\titlespacing*{\section}{0pt}{12pt}{4pt}
\pagestyle{plain}
\newcommand{\plainwords}[1]{\par\medskip\noindent\fcolorbox{black!25}{black!4}{\parbox{\dimexpr\linewidth-2\fboxsep-2\fboxrule}{\textbf{In plain words.} #1}}\par\medskip}
\newcommand{\titleblock}[3]{%
  {\LARGE\bfseries #1\par}\vspace{4pt}
  {\large #2\par}\vspace{2pt}
  {\small\color{black!70} Shakil A.\ Rafi, Ph.D.\ \textbar\ University of Arkansas \textbar\ #3\par}\vspace{2pt}
  \hrule\vspace{6pt}}

\begin{document}
\titleblock{nnR: Neural Networks Made Algebraic}{An R package for building neural networks by operating on them}{Software write-up, November 2026}
\textit{On CRAN (\href{https://cran.r-project.org/package=nnR}{cran.r-project.org/package=nnR}). Source: \href{https://github.com/2shakilrafi/nnR}{github.com/2shakilrafi/nnR}. License GPL-3. Co-authors: Joshua Lee Padgett and Ukash Nakarmi.}

\plainwords{Usually you pick a network shape and train it. Here, I treat a network as an object you can add, stack, compose, and multiply, then build networks that approximate functions like $x^2$, $e^x$, and $\sin x$ directly, with no training at all.}

\section{Introduction}
A feed-forward network can be written as an ordered tuple of weight and bias pairs, $\nu=((W_1,b_1),\dots,(W_L,b_L))$. Writing networks this way lets us define operations on them: composition, stacking, sums, scalar multiples, and networks for squaring and for multiplying two numbers. The framework builds on work by Grohs, Jentzen and collaborators, and our paper, \textit{Towards an Algebraic Framework for Approximating Functions Using Neural Network Polynomials} (Rafi, Padgett \& Nakarmi, 2024, \href{https://doi.org/10.48550/arXiv.2402.01058}{arXiv:2402.01058}), extends it to powers, polynomials, exponentials, sines, cosines, and a one-dimensional trapezoidal rule.

\section{Why a new package}
Common deep-learning libraries assume you define layers and then learn the parameters by backpropagation. This framework works the other way around: the weights and biases are fixed by construction, and some networks (the power network, for example) are defined entirely by their accuracy parameters. That does not fit a fixed layer structure well, so the operations needed their own implementation.

\section{What is implemented}
\begin{itemize}
\item Composition, stacking, sums, scalar multiplication, and instantiation of a network on an input.
\item Squaring, products, and powers; neural-network exponentials, sines, and cosines.
\item A one-dimensional trapezoidal rule, norms, maxima, and a maximum-convolution-type approximation for continuous functions of one variable.
\end{itemize}
Each function is tied to a definition in the paper, so the code can be read next to the mathematics. The package ships a vignette and a \texttt{testthat} test suite covering the operations.

\section{Status and discussion}
The package has been on CRAN since February 2024, and a newer release followed in September 2026. I describe it as a pedagogical proof of concept. It is not tuned for speed, and anyone needing performance should benchmark it first. The value I see is in making a piece of approximation theory executable and checkable, so students and researchers can test a construction against the definition on a real input instead of taking a proof on faith.

\textbf{Contact.} sarafi@uark.edu
\end{document}
